% source: RH-Framework/theorems/T01_E_NATIVE_REFINEMENT_COUNT_MEASURE.md
\hypertarget{t01-e1e2e3-native-refinement-count-content-and-simple-integrals}{%
\chapter{T01-E1/E2/E3 --- Native Refinement-Count Content and Simple Integrals}\label{t01-e1e2e3-native-refinement-count-content-and-simple-integrals}}

\textbf{Classification:} \texttt{PROVED}

\textbf{Foundational origin:} \texttt{NATIVE\_DERIVED}

\textbf{Continuous native measure classification:} \texttt{OPEN}

This theorem constructs a finite refinement content and its symmetric-difference
completion before introducing Haar measure, Lebesgue measure, a Jacobian, or an
ordinary \(L^p\) space.

It does \textbf{not} yet prove a sigma-additive measure on the actual continuous UGD
carrier. The completed object below is a complete finite-content algebra. The
continuous-function and actual-UGD adapters remain separate gates.

\begin{center}\rule{0.5\linewidth}{0.5pt}\end{center}

\hypertarget{native-refinement-grids}{%
\section{1. Native refinement grids}\label{native-refinement-grids}}

A refinement grid is a triple

\[
R=(n_s,n_t,n_\alpha),
\qquad
n_s,n_t,n_\alpha\in\mathbb N_{>0}.
\tag{1.1}
\]

The three entries count subdivisions in:

\begin{verbatim}
scale-refinement direction;
lifted phase-memory direction;
circular seam direction.
\end{verbatim}

At this stage these are native partition channels. They are not assumed to be
ordinary real coordinates.

A cell is indexed by

\[
c=(j_s,j_t,j_\alpha),
\qquad
j_s,j_t\in\mathbb Z,
\quad
j_\alpha\in\mathbb Z/n_\alpha\mathbb Z.
\tag{1.2}
\]

Define the reciprocal refinement content of one cell by

\[
\boxed{
\nu_R(c)=\frac{1}{n_sn_tn_\alpha}.
}
\tag{1.3}
\]

The value is a positive cut rational. No real length, volume form, or external
measure is primitive.

\begin{center}\rule{0.5\linewidth}{0.5pt}\end{center}

\hypertarget{finite-native-cylinders}{%
\section{2. Finite native cylinders}\label{finite-native-cylinders}}

A finite native cylinder is a finite set of cells on one grid:

\[
A=(R,C),
\qquad
C\subset
\mathbb Z\times\mathbb Z\times
\mathbb Z/n_\alpha\mathbb Z.
\tag{2.1}
\]

Its content is

\[
\boxed{
\nu(A)=\frac{|C|}{n_sn_tn_\alpha}.
}
\tag{2.2}
\]

Repeated cell indices are identified before counting. Thus content depends on
the represented finite cylinder, not on a list presentation.

\begin{center}\rule{0.5\linewidth}{0.5pt}\end{center}

\hypertarget{refinement}{%
\section{3. Refinement}\label{refinement}}

Let

\[
R'=(n_s',n_t',n_\alpha')
\]

refine \(R\), meaning each denominator of \(R'\) is a positive integer multiple
of the corresponding denominator of \(R\). Put

\[
q_s=\frac{n_s'}{n_s},
\qquad
q_t=\frac{n_t'}{n_t},
\qquad
q_\alpha=\frac{n_\alpha'}{n_\alpha}.
\tag{3.1}
\]

Each coarse cell is replaced by exactly

\[
q_sq_tq_\alpha
\tag{3.2}
\]

fine cells.

\hypertarget{theorem-3.1-refinement-invariance}{%
\subsection{Theorem 3.1 --- Refinement invariance}\label{theorem-3.1-refinement-invariance}}

For every finite cylinder \(A\),

\[
\boxed{
\nu(\operatorname{Ref}_{R'}A)=\nu(A).
}
\tag{3.3}
\]

\hypertarget{proof}{%
\subsection{Proof}\label{proof}}

If \(A\) contains \(m\) coarse cells, its refinement contains
\(mq_sq_tq_\alpha\) fine cells. Therefore

\[
\begin{aligned}
\nu(\operatorname{Ref}_{R'}A)
&=
\frac{mq_sq_tq_\alpha}
{n_s'n_t'n_\alpha'}\\
&=
\frac{m}{n_sn_tn_\alpha}
=
u(A).
\end{aligned}
\]

QED.

\begin{center}\rule{0.5\linewidth}{0.5pt}\end{center}

\hypertarget{common-refinement-and-boolean-operations}{%
\section{4. Common refinement and Boolean operations}\label{common-refinement-and-boolean-operations}}

For grids \(R\) and \(S\), define their common refinement coordinatewise by
least common multiples:

\[
R\vee S
=
(\operatorname{lcm}(n_s,m_s),
 \operatorname{lcm}(n_t,m_t),
 \operatorname{lcm}(n_\alpha,m_\alpha)).
\tag{4.1}
\]

Refine both cylinders to \(R\vee S\), then define union, intersection,
difference, and symmetric difference by ordinary finite-cell operations.
Refinement invariance shows that the resulting content is independent of the
chosen common refinement.

\hypertarget{theorem-4.1-finite-content-identity}{%
\subsection{Theorem 4.1 --- Finite-content identity}\label{theorem-4.1-finite-content-identity}}

For finite cylinders \(A,B\),

\[
\boxed{
\nu(A\cup B)+\nu(A\cap B)=\nu(A)+\nu(B).
}
\tag{4.2}
\]

In particular, if \(A\cap B=\varnothing\), then

\[
\nu(A\cup B)=\nu(A)+\nu(B).
\tag{4.3}
\]

\hypertarget{proof-1}{%
\subsection{Proof}\label{proof-1}}

Move to a common refinement. Every cell belongs to exactly one of

\[
A\setminus B,
\quad
B\setminus A,
\quad
A\cap B.
\]

Count the three disjoint finite sets and multiply by the common cell content.
QED.

The symmetric-difference identity is

\[
\boxed{
\nu(A\triangle B)
=
u(A)+\nu(B)-2\nu(A\cap B).
}
\tag{4.4}
\]

\begin{center}\rule{0.5\linewidth}{0.5pt}\end{center}

\hypertarget{native-translation-invariance}{%
\section{5. Native translation invariance}\label{native-translation-invariance}}

For integer shifts

\[
h=(h_s,h_t,h_\alpha),
\]

define

\[
T_h(j_s,j_t,j_\alpha)
=
(j_s+h_s,j_t+h_t,j_\alpha+h_\alpha),
\tag{5.1}
\]

with the seam coordinate reduced modulo \(n_\alpha\).

\hypertarget{theorem-5.1}{%
\subsection{Theorem 5.1}\label{theorem-5.1}}

\[
\boxed{
\nu(T_hA)=\nu(A).
}
\tag{5.2}
\]

\hypertarget{proof-2}{%
\subsection{Proof}\label{proof-2}}

Translation is a bijection of the cell index set on a fixed grid. It preserves
the number of cells and therefore preserves reciprocal-count content. QED.

This is native translation invariance on the refinement atlas. It is not yet a
Haar theorem on a locally compact group.

\begin{center}\rule{0.5\linewidth}{0.5pt}\end{center}

\hypertarget{symmetric-difference-gauge}{%
\section{6. Symmetric-difference gauge}\label{symmetric-difference-gauge}}

Define

\[
\boxed{
\delta(A,B)=\nu(A\triangle B).
}
\tag{6.1}
\]

After passing to common refinements, finite-set identities give

\[
\delta(A,B)\ge0,
\qquad
\delta(A,B)=\delta(B,A),
\tag{6.2}
\]

and

\[
A\triangle C
\subseteq
(A\triangle B)\cup(B\triangle C).
\tag{6.3}
\]

Consequently

\[
\boxed{
\delta(A,C)
\le
\delta(A,B)+\delta(B,C).
}
\tag{6.4}
\]

Two cylinder presentations have zero distance precisely when their common
refinements contain the same cells.

\hypertarget{lemma-6.1-content-continuity}{%
\subsection{Lemma 6.1 --- Content continuity}\label{lemma-6.1-content-continuity}}

\[
\boxed{
|\nu(A)-\nu(B)|\le\delta(A,B).
}
\tag{6.5}
\]

\hypertarget{proof-3}{%
\subsection{Proof}\label{proof-3}}

On a common refinement,

\[
\nu(A)-\nu(B)
=
u(A\setminus B)-\nu(B\setminus A).
\]

Taking the ordered rational absolute value yields

\[
|\nu(A)-\nu(B)|
\le
\nu(A\setminus B)+\nu(B\setminus A)
=
u(A\triangle B).
\]

QED.

\begin{center}\rule{0.5\linewidth}{0.5pt}\end{center}

\hypertarget{native-content-algebra-completion}{%
\section{7. Native content-algebra completion}\label{native-content-algebra-completion}}

A sequence \((A_n)\) of finite native cylinders is \textbf{refinement-content Cauchy}
when, for every positive cut rational \(\varepsilon\), there exists \(N\) such
that

\[
\delta(A_m,A_n)<\varepsilon
\qquad(m,n\ge N).
\tag{7.1}
\]

Two such sequences are equivalent when

\[
\delta(A_n,B_n)\longrightarrow0
\tag{7.2}
\]

in rational-threshold language.

Define

\[
\boxed{
\mathfrak M_\Sigma^{\mathrm{ref}}
=
\{\text{refinement-content Cauchy presentations}\}/\sim_\delta.
}
\tag{7.3}
\]

\hypertarget{theorem-7.1-complete-native-content-algebra}{%
\subsection{Theorem 7.1 --- Complete native content algebra}\label{theorem-7.1-complete-native-content-algebra}}

Union, intersection, difference, symmetric difference, and translation extend
to \(\mathfrak M_\Sigma^{\mathrm{ref}}\). The content extends uniquely by

\[
\boxed{
\widehat\nu([A_n])
=\lim_n\nu(A_n),
}
\tag{7.4}
\]

where the limit is taken in the cut-rational completion.

The extended content is nonnegative, finitely additive on disjoint completed
classes, translation invariant, and satisfies

\[
|\widehat\nu(X)-\widehat\nu(Y)|
\le
\widehat\delta(X,Y).
\tag{7.5}
\]

\hypertarget{proof-4}{%
\subsection{Proof}\label{proof-4}}

By \texttt{(6.5)}, \((\nu(A_n))\) is cut-Cauchy whenever \((A_n)\) is
\(\delta\)-Cauchy. Equation \texttt{(6.5)} also proves independence of representatives.

Finite-set inclusions give the Lipschitz estimates

\[
\delta(A\cup C,B\cup D)
\le
\delta(A,B)+\delta(C,D),
\tag{7.6}
\]

and analogous estimates for intersection, difference, and symmetric difference.
Thus Boolean-ring operations descend to Cauchy classes. Translation is an
isometry by \texttt{(5.2)}.

Finite additivity is inherited by approximating disjoint completed classes with
finite cylinders whose overlaps tend to zero in content. Completeness follows
from the Cauchy-presentation diagonal construction using \(\delta\)-entourages.
QED.

This theorem proves a complete finite-content algebra. It does not yet claim a
sigma-algebra or countable additivity for arbitrary completed families.

\begin{center}\rule{0.5\linewidth}{0.5pt}\end{center}

\hypertarget{nontrivial-cauchy-cylinder}{%
\section{8. Nontrivial Cauchy cylinder}\label{nontrivial-cauchy-cylinder}}

Let \(R_n=(2^n,1,1)\), and let \(S_n\) contain the first
\(2^n-1\) scale cells of one unit aperture. Then

\[
\nu(S_n)=1-2^{-n}.
\tag{8.1}
\]

For \(m>n\),

\[
\delta(S_n,S_m)
=2^{-n}-2^{-m}
<2^{-n}.
\tag{8.2}
\]

Thus \((S_n)\) is refinement-content Cauchy and represents a completed cylinder
of content one. The object is constructed through cell refinement, not by first
naming a real interval.

\begin{center}\rule{0.5\linewidth}{0.5pt}\end{center}

\hypertarget{native-simple-functions}{%
\section{9. Native simple functions}\label{native-simple-functions}}

A finite native simple function on a grid is a finite assignment

\[
f:C_R\to\mathbb K_\Sigma^0.
\tag{9.1}
\]

Zero values may be omitted. Refinement copies a coarse value to all fine cells
inside the coarse cell.

Define the tail integral

\[
\boxed{
\mathcal I_{\mathrm{tail}}(f)
=
\frac{1}{n_sn_tn_\alpha}
\sum_{c}D_\Sigma(f(c)),
}
\tag{9.2}
\]

and the recognition integral

\[
\boxed{
\mathcal I_{\mathrm{rec}}(f)
=
\frac{1}{n_sn_tn_\alpha}
\sum_c N_\Sigma(f(c)).
}
\tag{9.3}
\]

\hypertarget{theorem-9.1-refinement-invariance}{%
\subsection{Theorem 9.1 --- Refinement invariance}\label{theorem-9.1-refinement-invariance}}

For every refinement \(R'\succ R\),

\[
\mathcal I_{\mathrm{tail}}(\operatorname{Ref}_{R'}f)
=
\mathcal I_{\mathrm{tail}}(f),
\tag{9.4}
\]

\[
\mathcal I_{\mathrm{rec}}(\operatorname{Ref}_{R'}f)
=
\mathcal I_{\mathrm{rec}}(f).
\tag{9.5}
\]

\hypertarget{proof-5}{%
\subsection{Proof}\label{proof-5}}

Each value is copied to exactly \(q_sq_tq_\alpha\) fine cells, while fine cell
content is divided by the same factor. QED.

\hypertarget{theorem-9.2-translation-invariance}{%
\subsection{Theorem 9.2 --- Translation invariance}\label{theorem-9.2-translation-invariance}}

For every integer cell translation,

\[
\mathcal I_{\mathrm{tail}}(T_hf)
=
\mathcal I_{\mathrm{tail}}(f),
\qquad
\mathcal I_{\mathrm{rec}}(T_hf)
=
\mathcal I_{\mathrm{rec}}(f).
\tag{9.6}
\]

The proof is finite reindexing.

These integrals are native finite sums. Their Cauchy completions into general
integrable native functions remain T01-E4.

\begin{center}\rule{0.5\linewidth}{0.5pt}\end{center}

\hypertarget{foundational-exclusion-boundary}{%
\section{10. Foundational exclusion boundary}\label{foundational-exclusion-boundary}}

The construction above does not use as definitions or theorem sources:

\begin{verbatim}
Haar measure;
Lebesgue measure;
Caratheodory extension;
real-coordinate volume;
Jacobian density;
ordinary Lp spaces;
floating-point arithmetic.
\end{verbatim}

The implementation core uses exact integer counts and \texttt{Fraction} coordinates.
The Python type is a model of rational refinement content, not the mathematical
definition.

\begin{center}\rule{0.5\linewidth}{0.5pt}\end{center}

\hypertarget{later-coordinate-shadow}{%
\section{11. Later coordinate shadow}\label{later-coordinate-shadow}}

After an actual UGD refinement adapter is proved, one may attempt to identify the
three native partition channels with the normalized coordinates

\[
(u,t,\alpha).
\]

Only then may one prove that \(\widehat\nu\) corresponds to flat coordinate
content and subsequently that the original \((\phi,\sigma,k)\) chart has density
proportional to \(\sigma^{-2}\).

That future statement would be a coordinate-shadow theorem. The density may not
be used to generate the native content retroactively.

\begin{center}\rule{0.5\linewidth}{0.5pt}\end{center}

\hypertarget{exact-certificate}{%
\section{12. Exact certificate}\label{exact-certificate}}

The expected status is

\begin{verbatim}
PASS_NATIVE_REFINEMENT_COUNT_CONTENT_AND_SIMPLE_INTEGRALS
\end{verbatim}

The certificate checks exact refinement invariance, Boolean identities,
translation invariance, symmetric-difference continuity, a nontrivial Cauchy
slab, and refinement/translation invariance of both native simple integrals.

\begin{center}\rule{0.5\linewidth}{0.5pt}\end{center}

\hypertarget{remaining-gates}{%
\section{13. Remaining gates}\label{remaining-gates}}

\begin{verbatim}
T01-E1 NATIVE REFINEMENT CONTENT                  PROVED
T01-E2 COMPLETE NATIVE CONTENT ALGEBRA            PROVED
T01-E3 NATIVE SIMPLE TAIL/RECOGNITION INTEGRALS   PROVED

T01-E4 SIMPLE-FUNCTION CAUCHY COMPLETIONS          OPEN
T01-E5 ACTUAL UGD REFINEMENT-CHANNEL ADAPTER       OPEN
T01-E6 HAAR/JACOBIAN COORDINATE SHADOW             OPEN
ACTUAL MP PATH-ATLAS ADAPTER                       OPEN
NATIVE STATE-PACKET GENERATION                     OPEN
TARGET FAITHFULNESS                                OPEN
NATIVE SHUNYA                                      OPEN
RIEMANN HYPOTHESIS                                 OPEN
\end{verbatim}

\begin{flushright}\small\texttt{RH-Framework/theorems/T01\_E\_NATIVE\_REFINEMENT\_COUNT\_MEASURE.md}\end{flushright}
